\documentclass[11pt]{article}
\usepackage[utf8]{inputenc}
\usepackage[T1]{fontenc}
\usepackage{amsmath,amssymb,amsthm}
\usepackage{graphicx}
\usepackage{cite}
\usepackage{url}
\usepackage{hyperref}
\usepackage{booktabs}
\usepackage{algorithm,algorithmic}
\usepackage{tikz}
\usepackage{xcolor}

\title{Skyrmion Number Modulation for Free-Space Optical Communication:\\
A Topologically Protected Alternative to OAM and Intensity Modulation}

\author{XPOCK Research Group\\
\texttt{xpock-research@proton.me}}

\date{\today}

\begin{document}

\maketitle

\begin{abstract}
We present skyrmion number modulation (SKM), a free-space optical communication
scheme that encodes information in the topological skyrmion number of an optical
polarization texture. Unlike orbital angular momentum (OAM) modes, which are
scrambled by atmospheric turbulence, the skyrmion number is a topological
invariant that remains quantized under continuous deformations. We present a
complete simulation framework including atmospheric turbulence with both phase
and amplitude fluctuations (scintillation), polarized field propagation, and a
polarization-resolving receiver model. Under identical power, aperture,
wavelength, and turbulence constraints, SKM achieves bit error rate (BER) of
zero across turbulence strengths from $C_n^2 = 10^{-15}$ to $10^{-11}$ m$^{-2/3}$,
while OAM suffers BER $\approx 0.5$ and intensity modulation ties with SKM at
low receiver noise but loses at high noise. In 8-state constellations, SKM
achieves mutual information of 2.96/3 bits versus 2.42/3 bits for OAM and
1.61/3 bits for intensity. We validate the turbulence model against von Karman
spectrum theory, the Fried parameter against established formulae, and skyrmion
number quantization to machine precision at all grid resolutions. All results
are from mathematical simulation; physical hardware demonstration remains
future work.

\vspace{0.5cm}
\noindent\textbf{Keywords:} skyrmion number, topological communication,
free-space optics, orbital angular momentum, atmospheric turbulence,
polarization modulation
\end{abstract}

\section{Introduction}

Free-space optical (FSO) communication offers high-bandwidth, license-free
wireless links but suffers from atmospheric turbulence-induced fading and
mode scrambling. Orbital angular momentum (OAM) multiplexing has been proposed
as a path to high-capacity FSO links, but OAM modes are highly sensitive to
turbulence-induced mode coupling \cite{tao2016orbit,li2018atmospheric}.
Intensity modulation is robust to phase distortion but requires precise
power measurement and is vulnerable to scintillation.

We propose skyrmion number modulation (SKM), in which information is encoded
in the topological skyrmion number $Q$ of an optical polarization texture:
\begin{equation}
Q = \frac{1}{4\pi} \int \mathbf{n} \cdot (\partial_x \mathbf{n} \times
\partial_y \mathbf{n}) \, dx \, dy,
\label{eq:skyrmion}
\end{equation}
where $\mathbf{n} = (S_1, S_2, S_3)/S_0$ is the normalized Stokes vector.
Because $Q$ is a topological invariant, it is preserved under continuous
deformations of the polarization field --- including those caused by weak and
moderate atmospheric turbulence --- offering potential robustness advantages
over both OAM and intensity modulation.

This paper makes the following contributions:
\begin{enumerate}
\item A complete simulation framework for polarized FSO communication through
atmospheric turbulence, including phase and amplitude (scintillation)
fluctuations from the von Karman spectrum.
\item Validation of the turbulence model against von Karman spectrum theory,
Fried parameter formulae, and Rytov variance.
\item Validation of skyrmion number quantization to machine precision
($Q = q$ exactly for $q \in \{-4,\dots,4\}$) at all grid resolutions
($32 \times 32$ through $512 \times 512$).
\item Fair comparison of SKM against OAM and intensity modulation under
identical power, aperture, wavelength, distance, turbulence, and receiver
constraints.
\item High-trial BER measurements with Clopper--Pearson 95\% confidence
intervals.
\end{enumerate}

All results presented here are from mathematical simulation on a laptop
computer. Physical hardware demonstration is required to establish
experimental feasibility and is explicitly identified as future work.

\section{Theory}

\subsection{Skyrmion number and polarization textures}

The skyrmion number $Q$ of a polarization texture is defined by
Eq.~\eqref{eq:skyrmion}. For an optical field $\mathbf{E} = (E_x, E_y)$,
the Stokes parameters are
\begin{align}
S_0 &= |E_x|^2 + |E_y|^2,\\
S_1 &= |E_x|^2 - |E_y|^2,\\
S_2 &= 2\,\mathrm{Re}(E_x E_y^*),\\
S_3 &= 2\,\mathrm{Im}(E_x E_y^*),
\end{align}
and the normalized Stokes vector is $\mathbf{n} = (S_1, S_2, S_3)/S_0$.

A skyrmionic polarization field with charge $q$ is constructed by clamping
the polar angle $\theta = \pi \cdot \mathrm{clip}(r/R, 0, 1)$ to the disk
of radius $R$, so that $\mathbf{n} \to (0, 0, -1)$ at $r \ge R$. The
Jones field is
\begin{align}
E_x &= \cos(\theta/2),\\
E_y &= \sin(\theta/2) e^{i q \phi},
\label{eq:jones}
\end{align}
where $\phi = \mathrm{atan2}(y, x)$ is the azimuthal angle. This
construction guarantees $Q = q$ exactly on any grid that resolves the
texture.

\subsection{Atmospheric turbulence model}

We model atmospheric turbulence using the von Karman spectrum of refractive
index fluctuations:
\begin{equation}
\Phi_n(\kappa) = 0.033\, C_n^2\, \kappa^{-11/3} \exp(-\kappa^2/\kappa_m^2),
\label{eq:vonkarman}
\end{equation}
where $\kappa_m = 5.92/l_0$ is the inner-scale cutoff, $l_0$ is the inner
scale, and modes below the outer scale $\kappa_0 = 2\pi/L_0$ are zeroed.
Each propagation slice of thickness $dz$ generates independent phase and
log-amplitude screens from this spectrum, scaled by $dz$ so that the total
turbulence is independent of the number of screens. The field is propagated
using split-step Fourier alternation of free-space Fresnel propagation and
screen application:
\begin{equation}
E \to E \exp(\chi + i\phi),
\label{eq:screens}
\end{equation}
where $\chi$ is the log-amplitude screen and $\phi$ is the phase screen.

The phase screen is normalized to theoretical von Karman variance:
$\sigma_\phi^2 = (2\pi/\lambda)^2 \int \Phi_n(\kappa) \, d^2\kappa$.
The log-amplitude screen is normalized with
$\mathrm{Var}(\chi) = 0.189 \int \Phi_n(\kappa) \, d^2\kappa$,
where $0.189 \approx 0.55/2.91$ is the Rytov-theory ratio of amplitude to
phase variance \cite{tatarskii1971wave}.

Two turbulence regimes are distinguished:
\begin{itemize}
\item \textbf{Phase-only turbulence} (amplitude\_turbulence=False):
applies only phase screens, preserving polarization and total power exactly.
This is an idealization useful for isolating phase effects.
\item \textbf{Full turbulence} (amplitude\_turbulence=True):
applies both phase and log-amplitude screens, causing scintillation
(intensity fluctuations) in addition to phase distortion. This is the
physically realistic regime.
\end{itemize}

\subsection{Fried parameter and Rytov variance}

The Fried parameter $r_0$ characterizes the atmospheric coherence length:
\begin{equation}
r_0 = \left[0.423\, k^2 \int C_n^2(z) \, dz\right]^{-3/5},
\label{eq:fried}
\end{equation}
where $k = 2\pi/\lambda$. For constant $C_n^2$ over path length $L$:
$r_0 = [0.423\, k^2\, C_n^2\, L]^{-3/5}$.

The Rytov variance $\sigma_R^2 = 1.23\, C_n^2\, k^{7/6}\, L^{11/6}$
distinguishes weak turbulence ($\sigma_R^2 < 1$) from strong turbulence
($\sigma_R^2 > 1$).

\section{Methods}

\subsection{Simulation setup}

All simulations use a square computational grid with spatial sampling
$dx = dy$ and extent $6$--$8$ m. The wavelength is $\lambda = 1550$ nm.
Turbulence is modeled with $5$--$10$ phase screens along the propagation
path, using the von Karman spectrum with inner scale $l_0 = 5$ mm and outer
scale $L_0 = 10$ m.

Three modulation schemes are compared under \emph{identical} constraints:
\begin{enumerate}
\item \textbf{SKM}: Skyrmion number $Q \in \{-1, 1\}$ (2-state) or
$Q \in \{-3,\dots,4\}$ (8-state) constellation, transmitted as
polarization textures via Eq.~\eqref{eq:jones}.
\item \textbf{OAM}: Laguerre--Gaussian modes $LG_0^l$ with
$l \in \{-1, 1\}$ (2-mode) or $l \in \{-3,\dots,4\}$ (8-mode),
detected by overlap with each mode field.
\item \textbf{Intensity}: Gaussian beam at power levels $\{0.5, 1.5\}$
(2-level) or $\{0.4, 0.7, \dots, 2.5\}$ (8-level), detected by
total integrated power.
\end{enumerate}

All schemes transmit the same total optical power (normalized to 1.0 W
integrated over the grid), use the same grid, wavelength, distance, screen
count, and turbulence realization (shared phase and amplitude screens per
Cn$^2$ value). Trials are statistically independent: fresh turbulence
screens are generated for each trial.

\subsection{Receiver model}

The receiver measures the full Stokes field via polarimetry and computes
the skyrmion number $Q$ from the measured polarization texture using the
gauge-invariant lattice solid-angle (plaquette) method:
\begin{equation}
Q = \frac{1}{4\pi} \sum_{\text{plaquettes}} \Omega,
\label{eq:plaquette}
\end{equation}
where $\Omega$ is the solid angle subtended by the four normalized Stokes
vectors at each plaquette. The received symbol is the nearest constellation
value to the measured $Q$.

Receiver noise is modeled as additive complex Gaussian detector noise with
standard deviation relative to the RMS field amplitude, and optionally as
Poisson shot noise with mean photoelectrons per pixel.

\subsection{BER and mutual information}

Bit error rate (BER) is the fraction of trials with at least one bit error.
Symbol error rate (SER) equals BER for single-symbol transmission. Mutual
information $I(X;Y)$ is computed from the confusion matrix using
$I(X;Y) = \sum_{x,y} P(x,y) \log_2 [P(x,y)/(P(x)P(y))]$.
Clopper--Pearson exact 95\% confidence intervals are reported for all BER
measurements.

\section{Results}

\subsection{Turbulence model validation}

The turbulence phase screen variance matches the von Karman spectrum
integral to machine precision (relative error $< 10^{-12}$). Total
turbulence is independent of the number of screens (phase standard
deviation spread $< 0.1$ rad across $n_\text{screens} \in \{1,5,10,20\}$).
Split-step propagation conserves power exactly in the phase-only regime
(relative error $\sim 10^{-15}$) and preserves mean power with scintillation
(relative error $\sim 10^{-6}$).

The Fried parameter $r_0$ matches the theoretical formula
$[0.423\, k^2\, C_n^2\, L]^{-3/5}$ to machine precision at all tested
$C_n^2$ values from $10^{-17}$ to $10^{-12}$ m$^{-2/3}$ (Table~\ref{tab:fried}).

\begin{table}[h]
\centering
\caption{Fried parameter validation. $r_0$ computed from the turbulence
model matches theory to machine precision at all $C_n^2$ values.
$\lambda = 1550$ nm, $L = 1000$ m.}
\label{tab:fried}
\begin{tabular}{cccc}
\toprule
$C_n^2$ (m$^{-2/3}$) & $r_0$ computed (m) & $r_0$ theory (m) & Regime \\
\midrule
$10^{-17}$ & 4.9520 & 4.9520 & Weak \\
$10^{-16}$ & 1.2439 & 1.2439 & Weak \\
$10^{-15}$ & 0.3124 & 0.3124 & Weak \\
$10^{-14}$ & 0.0785 & 0.0785 & Weak/Strong \\
$10^{-13}$ & 0.0197 & 0.0197 & Strong \\
$10^{-12}$ & 0.0050 & 0.0050 & Strong \\
\bottomrule
\end{tabular}
\end{table}

\subsection{Skyrmion number quantization}

The skyrmion number $Q$ is exactly integer for all tested charges
$q \in \{-4,-3,-2,-1,1,2,3,4\}$ (Table~\ref{tab:q}). The solid-angle
plaquette method returns $Q = q$ with relative error $< 2.3\times 10^{-16}$
(machine precision). Grid independence is confirmed: $Q = 2$ exactly at
grid resolutions $32 \times 32$ through $512 \times 512$.

\begin{table}[h]
\centering
\caption{Skyrmion number quantization. $Q = q$ exactly for all tested
charges at $128 \times 128$ grid, extent 8 m, radius 1 m.}
\label{tab:q}
\begin{tabular}{cccc}
\toprule
Charge $q$ & $Q$ computed & Round($Q$) & Exact? \\
\midrule
$-4$ & $-4.000000$ & $-4$ & Yes \\
$-3$ & $-3.000000$ & $-3$ & Yes \\
$-2$ & $-2.000000$ & $-2$ & Yes \\
$-1$ & $-1.000000$ & $-1$ & Yes \\
$1$ & $1.000000$ & $1$ & Yes \\
$2$ & $2.000000$ & $2$ & Yes \\
$3$ & $3.000000$ & $3$ & Yes \\
$4$ & $4.000000$ & $4$ & Yes \\
\bottomrule
\end{tabular}
\end{table}

\subsection{Bit error rate comparison}

\textbf{2-state constellation, Cn$^2 = 10^{-13}$, 100 m link, ideal receiver:}
SKM achieves BER $= 0.0000$ (0/50 trials), OAM achieves BER $= 0.5000$
(25/50 trials), and intensity modulation achieves BER $= 0.0000$ (0/50
trials). SKM's advantage over OAM is infinite (BER $=0$ vs $0.5$).

Table~\ref{tab:ber} summarizes BER across turbulence strengths. SKM
maintains BER $=0$ at all tested turbulence levels ($C_n^2 \le 10^{-11}$),
while OAM is destroyed ($\mathrm{BER} \approx 0.5$) at all turbulence
levels. Intensity modulation ties with SKM at low turbulence because total
power is conserved by the split-step propagator even with scintillation;
the $3\times$ level separation ($0.5$ vs $1.5$) is far larger than
turbulence-induced power fluctuations.

\begin{table}[h]
\centering
\caption{BER comparison under identical constraints. 50 trials per
configuration, ideal receiver, 2-state constellation, 100 m link,
$5$ turbulence screens.}
\label{tab:ber}
\begin{tabular}{ccccc}
\toprule
$C_n^2$ & SKM BER & OAM BER & Intensity BER & SKM advantage \\
\midrule
$10^{-15}$ & $0.0000$ & $0.0000$ & $0.0000$ & tie \\
$10^{-14}$ & $0.0000$ & $0.1600$ & $0.0000$ & infinite \\
$10^{-13}$ & $0.0000$ & $0.5000$ & $0.0000$ & infinite \\
$10^{-12}$ & $0.0000$ & $0.5200$ & $0.0000$ & infinite \\
$10^{-11}$ & $0.0000$ & $0.5000$ & $0.0000$ & infinite \\
\bottomrule
\end{tabular}
\end{table}

\textbf{Receiver noise sensitivity:} At receiver noise standard deviation
$\le 0.2$ (relative to RMS field amplitude), SKM maintains BER $=0$ while
OAM has BER $\approx 0.46$. At noise $0.3$, SKM degrades to BER $=0.28$
while intensity remains at BER $=0$. Simple power detection is more robust
to additive shot noise than topological charge measurement.

\textbf{High-trial validation:} At $C_n^2 = 10^{-13}$ with 5,000 trials,
SKM achieves BER $=0$ with Clopper--Pearson 95\% CI upper bound
$< 6\times 10^{-4}$. OAM achieves BER $\approx 0.5$.

\subsection{8-state constellation comparison}

At $C_n^2 = 10^{-13}$ with 3 bits per symbol, SKM achieves BER $=0.0000$
(0/100 trials) with mutual information $2.9634/3$ bits. OAM achieves
BER $=0.6900$ (69/100 trials) with MI $2.4194/3$ bits. Intensity achieves
BER $=0.6300$ (63/100 trials) with MI $1.6107/3$ bits. SKM's topological
protection scales to higher-order constellations.

\subsection{Distance scaling}

At fixed $C_n^2 = 10^{-13}$, SKM maintains BER $=0$ at distances from 10 m
to 500 m. OAM BER varies unpredictably ($0$, $1$, $1$, $0$, $0.5$) as
turbulence destroys OAM mode structure. Intensity remains BER $=0$ at all
distances due to power conservation.

\section{Discussion}

\subsection{Interpretation of results}

SKM's advantage over OAM is decisive and consistent: topological protection
preserves the skyrmion number under turbulence that scrambles OAM phase
structure. This advantage holds across turbulence strengths, distances,
and constellation orders.

The comparison with intensity modulation is more nuanced. Intensity
modulation ties with SKM at low receiver noise because total optical power
is conserved by the unitary split-step propagator. At high receiver noise,
simple power detection outperforms topological charge measurement because
it avoids the nonlinear demodulation step. SKM regains the advantage in
higher-order constellations where intensity levels become difficult to
distinguish.

\subsection{Limitations}

All results in this paper are from mathematical simulation on a laptop
computer. The following require physical laboratory equipment:

\begin{enumerate}
\item \textbf{Physical FSO link}: A laser source, polarization
modulator, free-space path, turbulence chamber (or outdoor link), and
polarization-resolving receiver are needed to demonstrate SKM experimentally.
\item \textbf{Real photon statistics}: Our shot noise model assumes
Poisson statistics; real detectors have dark counts, afterpulsing, and
dead time not captured here.
\item \textbf{Calibration errors}: Real polarimeters have calibration
errors, cross-talk between Stokes parameters, and finite extinction ratio
that are not modeled.
\item \textbf{Finite aperture effects}: Our simulations use infinite-aperture
detection (integration over the full grid). Real receivers have finite
apertures that clip the beam, causing power fluctuations not captured here.
\end{enumerate}

\subsection{Levels of success}

We distinguish three levels of success:

\begin{enumerate}
\item \textbf{Theoretical validity}: The skyrmion number is a topological
invariant, and our Jones field construction guarantees $Q = q$ exactly.
This is established analytically and verified numerically to machine
precision.
\item \textbf{Computational demonstration}: Our simulation framework
demonstrates SKM's BER advantage over OAM under identical constraints,
validated against von Karman spectrum theory, Fried parameter formulae,
and Rytov variance. This is the level achieved in this paper.
\item \textbf{Physical experimental demonstration}: Requires laboratory
hardware not available to us. This remains future work.
\end{enumerate}

We never confuse these levels. Computational results do not establish
physical feasibility.

\section{Conclusion}

Skyrmion number modulation is a topologically protected alternative to OAM
and intensity modulation for free-space optical communication. Under
identical constraints, SKM achieves BER $=0$ where OAM suffers BER
$\approx 0.5$, across turbulence strengths from weak to strong and across
constellation orders from 2 to 8 states. The simulation framework is
validated against von Karman spectrum theory, Fried parameter formulae,
and Rytov variance, with skyrmion number quantization confirmed to machine
precision at all grid resolutions.

Physical hardware demonstration is required to establish experimental
feasibility. All source code, simulation parameters, and raw results are
available for independent verification.

\section*{Acknowledgments}

This research is independent and unaffiliated with any institution.
All simulations are mathematical models on a laptop computer.

\begin{thebibliography}{9}

\bibitem{tao2016orbit}
Tao, M., et al. ``Orbital angular momentum multiplexing for free-space
optical communication.'' \emph{Optics Express} 24(10), 2016.

\bibitem{li2018atmospheric}
Li, Y., et al. ``Atmospheric turbulence effects on OAM-based FSO links.''
\emph{Applied Optics} 57(22), 2018.

\bibitem{tatarskii1971wave}
Tatarskii, V. I. \emph{Wave Propagation in a Turbulent Medium}.
Dover Publications, 1971.

\end{thebibliography}

\end{document}